\section{Introduction}
\label{sec:intro}

Haematoxylin and eosin (H\&E) slides of the same tissue look different when they are digitised on
different whole-slide scanners. Light source, sensor and colour calibration all shift the pixel colour
distribution, and computational pathology models learn features that are partly entangled with it, so
accuracy can fall when a model trained on one scanner meets another~\citep{breen2024gans}. Stain
normalisation addresses this as a preprocessing step: it maps each input towards a target colour
distribution before any downstream model sees it, without retraining that model.

Classical normalisers~\citep{macenko2009normalisation,reinhard2001colortransfer} are fast and need no
training data, but they assume a linear stain model. GAN-based methods such as
StainGAN~\citep{shaban2019staingan} learn non-linear mappings but may alter tissue structure, and fast
learned colour maps such as StainNet~\citep{kang2021stainnet} and ParamNet~\citep{kang2024paramnet} are
usually trained on large collections. Latent diffusion models~\citep{rombach2022ldm} offer a strong image
prior and a mature adapter ecosystem: LoRA~\citep{hu2021lora} for data-efficient adaptation,
ControlNet~\citep{zhang2023controlnet} for spatial conditioning and LCM-LoRA~\citep{luo2023lcmlora} for
few-step inference. A similar combination has been used for domain-targeted style transfer in plant
imagery~\citep{hartley2024plant}, but not, to our knowledge, for scanner normalisation in histopathology.

This paper tests that combination under a deliberately strict few-shot budget: at most 50
coordinate-corresponding (not pixel-exact) crop pairs from a single MITOS-ATYPIA-14
slide~\citep{roux2014mitos}, evaluated on five held-out slide pairs that no model ever sees in training.
Our contributions are:
\begin{enumerate}
\item a six-rung ablation (A0--A5) that isolates what ControlNet, a colour LoRA, LCM-LoRA and a
histopathology warm-start each contribute;
\item a source-conditioned translator with deterministic DDIM inversion, the first configuration that
moves colour towards the target scanner on every held-out slide while also improving structure over the
earlier rungs;
\item an analysis of the VAE reconstruction ceiling, a post-hoc colour-residual fusion that sidesteps it,
and a decoder fine-tune that raises the ceiling mainly by smoothing;
\item a matched-budget comparison against three classical and three learned baselines, including a
structure-safety test with a pretrained nucleus segmenter; and
\item a portability test on SDXL~\citep{podell2023sdxl}.
\end{enumerate}
We report negative results as negative. In particular, two learned baselines trained on the same 50 pairs
outperform every diffusion configuration we built.
